\begin{textsample}{POS Dim 2 – 2019\_c5 – Score -1.00 – 2019\_c5/enade\_2019\_eng\_amb\_0cc6810176.txt}  \label{ex:f2_pos_337}
CAPÍTULO 5 DISTRIBUIÇÃO DOS CONCEITOS 5.1 PANORAMA NACIONAL DA DISTRIBUIÇÃO DOS CONCEITOS O Conceito Enade é um indicador de qualidade que avalia os cursos por intermédio dos desempenhos dos estudantes a partir dos resultados do Enade.
Os valores possíveis vão de 1 ( pior situação ) a 5 ( melhor situação ).
Os cursos com apenas um ou sem nenhum concluinte participante não obtêm o Conceito Enade, ficando Sem Conceito ( SC ).
Na Tabela 5.1, são apresentadas a quantidade e a distribuição de cursos de Engenharia Ambiental participantes do Enade/2019, por faixa de conceito e Grande Região.
Enfatiza -se, mais uma vez, que a diferença entre os cursos tabulados neste capítulo e no capítulo 2 corresponde aos cursos que não foram avaliados, ou seja, àqueles, em princípio, sem alunos concluintes que tivessem sido inscritos para a prova.
Observando -se os dados da Tabela 5.1, nota -se que, dos 292 cursos participantes, 116 ( 39,7 \% ) classificaram -se com conceito 3, o valor modal.
Esse foi também o conceito modal nas regiões Norte ( 40,0 \%, ao conceito 2 também foram atribuídos 40 \% dos cursos ), Nordeste ( 31,3 \% ), Sudeste ( 41,4 \% ) e Sul ( 50,0 \% ).
Os conceitos 2 e 4 foram o segundo mais frequente em nível nacional ( 22,3 \%, correspondendo a 65 cursos cada um ), e o conceito 5, o quarto ( 5,8 \%, correspondendo a 17 cursos ).
Houve, ainda, 14 cursos ( 4,8 \% ) que receberam conceito 1, e 15 cursos ( 5,1 \% ) ficaram Sem Conceito ( SC ).
A região Norte participou com 20 cursos ou 6,8 \% do total nacional.
Desses, oito cursos receberam os conceitos 2 e 3, o que equivale a 40,0 \% do total regional para cada conceito.
A três cursos ( 15,0 \% ) atribuiu -se conceito 1, e a um curso ( 5,0 \% ), conceito 4.
Nenhum curso ficou Sem Conceito ( SC ) ou recebeu conceito 5.
No Gráfico 5.1, é apresentada a distribuição cumulativa do Conceito Enade segundo Grande Região.
Em um gráfico de distribuição cumulativa, em cada valor de abscissa ( eixo x ), o valor apresentado no eixo das ordenadas ( eixo y ) é o valor acumulado da variável ( conceito Enade ).
Observando a linha vertical pontilhada, que ajuda a visualizar o percentual de conceitos acumulados até o valor 3, ou seja, conceitos 1, 2 e 3, pode -se notar que, por exemplo, a região Norte ( linha azul ) apresenta 95 \% dos cursos com conceito 3 ou abaixo deste valor.
A região Centro-Oeste ( linha laranja ), por outro lado, apresenta 65 \% dos cursos na mesma situação : com conceito 3 ou abaixo.
Destaca -se que em algumas regiões o conceito 5 pode não ter sido alcançado por nenhum curso, como ocorre na região Norte ( linha azul ), na qual todos os cursos têm conceito menor ou igual a 4, e sua poligonal alcança 100 \% no conceito 4.
Em linhas gerais, em um gráfico de distribuição cumulativa, poligonais mais à esquerda representam Grandes Regiões com pior distribuição, e poligonais mais à direita, Grandes Regiões com melhores conceitos.
A região Norte ( linha azul ) apresenta os cursos com A região Nordeste participou com 48 cursos ou 16,4 \% do total nacional.
Nessa região, 15 cursos ( 31,3 \% em termos regionais ) obtiveram conceito 3, o conceito modal para a região, como já comentado.
Os conceitos 2 e 4 foram atribuídos, respectivamente, a dez e a 12 cursos ( 20,8 \% e 25,0 \%, respectivamente ).
Ainda nessa região, quatro cursos ( 8,3 \% ) foram avaliados com o conceito 5, e três cursos ( 6,3 \% ) receberam conceito 1.
Quatro cursos ficaram Sem Conceito ( SC ).
Dos 140 cursos participantes da região Sudeste ( 47,9 \% do total nacional ), o conceito 3 foi o conceito modal, como já anteriormente assinalado, tendo sido obtido por 58 cursos ( 41,4 \% ).
O conceito 1 foi atribuído a cinco cursos ( 3,6 \% ), e o conceito 2, a 29 cursos ( 20,7 \% ).
Receberam os conceitos 4 e 5, respectivamente, 31 cursos ( 22,1 \% ) e oito cursos ( 5,7 \% ).
Nessa região, nove cursos ficaram Sem Conceito ( SC ).
Os 60 cursos da região Sul corresponderam a 20,5 \% do total nacional.
Como já indicado, houve predominância do conceito 3, atribuído a 50,0 \% dos cursos da região ( conceito modal ), ou seja, atribuído a 30 dos 60 cursos participantes na região Sul.
O conceito 1 foi atribuído a dois cursos ( 3,3 \% ), conceito 2, a nove cursos ( 15,0 \% ), conceito 4, a 15 cursos ( 25,0 \% ) e o conceito 5, a três cursos ( 5,0 \% ).
Um curso ( 1,7 \% ) ficou Sem Conceito.
Já dos 24 cursos participantes na região Centro-Oeste ( 8,2 \% do total nacional ), nove ( 37,5 \% em termos regionais ) receberam conceito 2, o conceito modal.
Seis cursos ( 25,0 \% ) obtiveram conceito 4, cinco cursos ( 20,8 \% ) obtiveram conceito 3, dois cursos ( 8,3 \% ) obtiveram conceito 5.
Um curso ficou Sem Conceito ( SC ) e um curso recebeu conceito 1.
No Gráfico 5.1, é apresentada a distribuição cumulativa do Conceito Enade segundo Grande Região.
Em um gráfico de distribuição cumulativa, em cada valor de abscissa ( eixo x ), o valor apresentado no eixo das ordenadas ( eixo y ) é o valor acumulado da variável ( conceito Enade ).
Observando a linha vertical pontilhada, que ajuda a visualizar o percentual de conceitos acumulados até o valor 3, ou seja, conceitos 1, 2 e 3, pode -se notar que, por exemplo, a região Norte ( linha azul ) apresenta 95 \% dos cursos com conceito 3 ou abaixo deste valor.
A região Centro-Oeste ( linha laranja ), por outro lado, apresenta 65 \% dos cursos na mesma situação : com conceito 3 ou abaixo.
Destaca -se que em algumas regiões o conceito 5 pode não ter sido alcançado por nenhum curso, como ocorre na região Norte ( linha azul ), na qual todos os cursos têm conceito menor ou igual a 4, e sua poligonal alcança 100 \% no conceito 4.
Em linhas gerais, em um gráfico de distribuição cumulativa, poligonais mais à esquerda representam Grandes Regiões com pior distribuição, e poligonais mais à direita, Grandes Regiões com melhores conceitos.
A região Norte ( linha azul ) apresenta os cursos com conceitos mais baixos, ao passo que as regiões Nordeste ( linha verde ) e Centro-Oeste ( linha laranja ) apresentam as distribuições com valores maiores.
Cursos em IES situadas nas regiões Sudeste ( linha vermelha ) e Sul ( linha rosa ) apresentam uma situação intermediária. 5.2 CONCEITOS POR CATEGORIA ADMINISTRATIVA, POR MODALIDADE DE ENSINO E GRANDE REGIÃO Na Tabela 5.2, é apresentada a distribuição dos cursos de Engenharia Ambiental participantes do Enade/2019, por Categoria Administrativa e Modalidade de ensino, de acordo com os conceitos por eles alcançados, segundo as Grandes Regiões brasileiras.
Dos 292 cursos participantes, 98 ( 33,6 \% ) eram ministrados em Instituições Públicas, e 194 ( 66,4 \% ), em Instituições Privadas.
Quanto à Modalidade de ensino, seis cursos eram a Distância, e 286, Presenciais.
De acordo com as informações da Tabela 5.2, em termos nacionais, entre Instituições Públicas e Privadas, observa -se que, dos 17 cursos avaliados com conceito 5, 16 eram oferecidos em IES Públicas, e um, em IES Privadas.
Dos 98 cursos participantes de IES Públicas, o conceito 4 foi o valor modal, atribuído a 43 cursos ( 43,9 \% ).
Dos demais cursos de Instituições Públicas participantes, três cursos ( 3,1 \% ) foram avaliados com conceito 1, seis cursos ( 6,1 \% ) foram avaliados com conceito 2, 27 cursos ( 27,6 \% ) foram avaliados com conceito 3 e, 16 cursos ( 16,3 \% ) receberam conceito 5.
Três cursos ficaram Sem Conceito ( SC ).
Na rede privada, o conceito modal foi 3, atribuído a 89 cursos, o correspondente a 45,9 \% dos 194 cursos da categoria.
Dos demais cursos participantes, 11 ( 5,7 \% ) receberam conceito 1, e 59 ( 30,4 \% ), conceito 2.
O conceito 4 foi atribuído a 22 cursos ( 11,3 \% ), e o conceito 5, a um curso ( 0,5 \% ).
Nesta Categoria Administrativa, 12 cursos ( 6,2 \% ) ficaram Sem Conceito ( SC ).
Dos 286 cursos participantes da modalidade presencial, o conceito 3 foi o valor modal, atribuído a 114 cursos ( 39,9 \% ).
Dos demais cursos nesta modalidade de ensino, 14 cursos ( 4,9 \% ) receberam conceito 1, 62 cursos ( 21,7 \% ) receberam conceito 2, 64 cursos ( 22,4 \% ), conceito 4, e 17 cursos ( 5,9 \% ), conceito 5.
Nesta modalidade de ensino, 15 cursos ( 5,2 \% ) ficaram Sem Conceito ( SC ).
Aos seis cursos a Distância informados, atribuíram -se os conceitos 2, 3 e 4.
Os conceitos foram atribuídos a três, dois e um cursos, respectivamente.
Nessa Modalidade de ensino, nenhum curso ficou Sem Conceito ( SC ).
Na análise por região, observa -se que, na região Norte, as Instituições Privadas participaram com 12 cursos ( 60,0 \% do total regional ), nenhum ficou Sem Conceito ( SC ), e nenhum recebeu os conceitos 4 ou 5.
O conceito modal para as Instituições Privadas na região foi 2, atribuído a sete cursos, correspondendo a 58,3 \% dessa categoria na região.
Dois cursos foram avaliados com conceito 1, três cursos, com conceito 3.
As Instituições Públicas participaram com oito cursos na região Norte ( 40,0 \% do total regional ).
Cinco cursos receberam conceito 3 ; os conceitos 1, 2 e 4 foram atribuídos a um curso cada.
Nenhum curso ficou Sem Conceito ( SC ) ou recebeu conceito 5, nessa combinação de categoria e região.
Todos os 20 cursos oferecidos na região Norte eram presenciais.
Na região Nordeste, a rede privada concentrou 27 dos 48 cursos participantes, o equivalente a 56,3 \% do total da região.
Dentre os cursos oferecidos por IES Privadas no Nordeste, 12 cursos foram avaliados com conceito 3, o conceito modal.
Dez cursos obtiveram conceito 2, dois cursos receberam conceito 1 e um curso ficou com conceito 5.
Nessa combinação de categoria e região, dois cursos ficaram Sem Conceito ( SC ) e nenhum recebeu conceito 4.
As Instituições Públicas da região Nordeste participaram com 19 cursos ( 43,2 \% do total da região ).
Desses, 12 obtiveram conceito 4, o valor modal, os conceitos 3 e 5 receberam três cursos cada.
Um curso recebeu conceito 1.
Dois cursos ficaram Sem Conceito ( SC ) e nenhum recebeu conceito 2.
Todos os cursos oferecidos na região Nordeste eram presenciais.
Na região Sudeste, a proporção de cursos da rede privada, 75,7 \%, corresponde a 106 dos 140 cursos participantes.
Nessa categoria e região, o conceito modal foi 3, atribuído a 53 cursos, e nove cursos ficaram Sem Conceito ( SC ).
Os demais foram avaliados com conceito 1 ( cinco cursos ), conceito 2 ( 26 cursos ) e conceito 4 ( 13 cursos ).
Nenhum curso recebeu conceito 5.
Dos 34 cursos oferecidos por Instituições Públicas na região Sudeste, o conceito modal foi 4, atribuído a 18 cursos.
Nessa região, oito cursos de IES Públicas receberam conceito 5, cinco cursos receberam conceito 3 e três receberam conceito 2.
Nenhum curso ficou Sem Conceito ( SC ) ou foi avaliado com conceito 1.
Também na região Sudeste, a maioria dos cursos foi oferecida na modalidade Presencial : 134 cursos.
Para esses, a moda foi o conceito 3, com 56 cursos.
Os demais cursos presenciais da região foram alocados nos conceitos 1 ( cinco cursos ), 2 ( 26 cursos ), 4 ( 30 cursos ) e 5 ( oito cursos ).
Ficaram Sem Conceito ( SC ) nove cursos.
Na região Sudeste, os seis cursos oferecidos na modalidade Educação a Distância obtiveram os conceitos 2, 3 e 4, respectivamente, três, dois e um curso, respectivamente.
As Instituições Privadas concentraram 38 dos 60 cursos participantes da região Sul, 63,3 \% do total regional.
Desses, 20 cursos obtiveram conceito 3, o conceito modal.
Nesta combinação de Categoria Administrativa e Grande Região, dois cursos receberam conceito 1, nove cursos receberam conceito 2, seis cursos receberam conceito 4.
Nenhum curso recebeu conceito 5, e um curso ficou Sem Conceito ( SC ).
As Instituições Públicas na região Sul participaram com 22 cursos ( 36,7 \% dos cursos da região ), dez cursos receberam conceito 3, nove cursos receberam conceito 4 e três receberam conceito 5.
Nenhum curso ficou Sem Conceito ( SC ) ou recebeu os demais conceitos.
Todos os cursos oferecidos na região Sul eram presenciais.
Na região Centro-Oeste, 11 dos 24 cursos participantes eram de Instituições Privadas ( 45,8 \% em termos regionais ).
Desses, sete receberam conceito 2, o conceito modal.
Dos demais cursos, um recebeu conceito 3, e outros três, conceito 4.
Quanto aos 13 cursos oferecidos por Instituições Públicas na região Centro-Oeste ( 54,2 \% do total regional ), quatro cursos apresentaram conceito 3, conceito modal, três receberam conceito 4.
Os conceitos 2 e 5 foram atribuídos a dois cursos cada.
Um curso ficou Sem Conceito ( SC ) e um curso recebeu conceito 1, nessa combinação de categoria e região.
Todos os cursos oferecidos na região Centro-Oeste eram presenciais.
No Gráfico 5.2, apresenta -se a distribuição do Conceito Enade, segundo a Categoria Administrativa da IES.
Os cursos em IES Públicas ( linha azul ) apresentam uma poligonal mais à direita do que os em IES Privadas ( linha verde ) e, consequentemente, uma distribuição de Conceitos Enade com valores maiores.
Como nenhum curso de IES Privada alcançou o conceito 5, a linha poligonal correspondente ( em verde ) alcança 100 \% no conceito 4. 5.3 CONCEITOS POR ORGANIZAÇÃO ACADÊMICA E GRANDE REGIÃO Na Tabela 5.3, encontra -se a distribuição dos conceitos atribuídos aos cursos participantes do Enade/2019, na Área de Engenharia Ambiental, por Organização Acadêmica, segundo as Grandes Regiões brasileiras.
Dos 292 cursos de Engenharia Ambiental participantes, 146 eram oferecidos em Universidades, 67 em Centros Universitários, 70 em Faculdades e 9 em CEFET/IF.
Esta distribuição corresponde a, respectivamente, 50,0 \%, 22,9 \%, 24,0 \% e 3,1 \% dos cursos.
De acordo com os dados apresentados, de todos os 17 cursos avaliados com conceito 5, 16 eram vinculados a Universidades.
Os cursos vinculados a esse tipo de Organização Acadêmica tiveram o conceito 3 como conceito modal, atribuído a 54 cursos ( 37,0 \% ).
Os demais cursos vinculados a Universidades receberam os conceitos 1 ( seis cursos ), 2 ( 17 cursos ), 4 ( 47 cursos ) e conceito 5 ( 16 cursos, como já mencionado ).
Ficaram Sem Conceito ( SC ), seis cursos.
Dos cursos em Centros Universitários, o conceito modal também foi 3, atribuído a 31 dos 67 cursos ( 46,3 \% ).
Cinco cursos vinculados a esse tipo de Organização Acadêmica ficaram Sem Conceito ( SC ), e os outros cursos receberam os conceitos 1 ( dois cursos ), 2 ( 23 cursos ), 4 ( cinco cursos ) e 5 ( um curso ).
Dos 70 cursos mantidos por Faculdades, 29 ( 41,4 \% ) receberam conceito 3, o conceito modal.
O segundo conceito mais frequentemente atribuído aos cursos oferecidos por essa categoria de IES foi o conceito 2, recebido por 25 cursos.
Os conceitos 1 e 4 foram atribuídos a seis cursos cada.
Quatro cursos ficaram Sem Conceito ( SC ) e nenhum recebeu conceito 5.
Dos nove cursos mantidos por CEFET/IF, sete ( 77,8 \% ) receberam conceito 4, o conceito modal.
O segundo conceito mais frequentemente atribuído aos cursos oferecidos por essa categoria de IES foi o conceito 3, recebido por dois cursos.
Os demais conceitos não foram atribuídos a nenhum curso.
Nenhum curso ficou Sem Conceito ( SC ).
Considerando -se, separadamente, as regiões brasileiras, verifica -se que, na região Norte, as Universidades participaram com 10 dos 20 cursos avaliados.
O conceito modal foi 3, atribuído a seis cursos.
Os conceitos 2 e 4 foram atribuídos a um curso cada.
Dois cursos receberam conceito 1.
Nenhum curso ficou Sem Conceito ( SC ) ou recebeu conceito 5.
Os Centros Universitários da região Norte foram representados por quatro cursos, os quais obtiveram conceitos 2 ( três cursos ) e 3 ( um curso ).
Nenhum curso ficou Sem Conceito ( SC ) ou recebeu os demais conceitos nesta combinação de categoria e região.
As Faculdades participaram com seis cursos na região Norte, os quais receberam os conceitos 2 ( quatro cursos, conceito modal ), 3 ( um curso ) e 1 ( um curso ).
Já os CEFET/IF não possuíam cursos avaliados na região Norte.
Na região Nordeste, as Universidades participaram com 23 dos 48 cursos da Área de Engenharia Ambiental oferecidos.
O conceito modal foi 4, com nove cursos.
Seguido do conceito 3, com seis cursos e do conceito 5, com três cursos.
Os conceitos 1 e 2 foram atribuídos a um curso cada.
Três cursos ficaram Sem Conceito ( SC ).
Os Centros Universitários contaram com 12 cursos participantes na região Nordeste, seis dos quais receberam o conceito modal 3.
Para os demais cursos, foram atribuídos os conceitos 2 ( cinco cursos ) e 5 ( um curso ).
Nenhum curso oferecido por Centros Universitários no Nordeste ficou Sem Conceito ( SC ) ou recebeu os conceitos 1 ou 4.
As Faculdades foram representadas por nove cursos na região Nordeste, quatro dos quais receberam conceito 2, o conceito modal.
Cada um dos conceitos 1 e 3 foi atribuído a dois cursos.
Um curso ficou Sem Conceito ( SC ).
Nenhum curso recebeu os demais conceitos.
Os CEFET/IF participaram com quatro cursos na região Nordeste, avaliados com conceitos 4 ( três cursos ) e 3 ( um curso ).
Nenhum curso ficou Sem Conceito ( SC ), ou recebeu os conceitos 1, 2 ou 5.
Na região Sudeste, as Universidades concentraram 58 dos 140 cursos de Engenharia Ambiental da região.
Dos cursos oferecidos em Universidades, nessa região, o conceito modal foi 3, atribuído a 21 cursos.
Os demais cursos receberam os conceitos 1 ( um curso ), 2 ( oito cursos ), 4 ( 19 cursos ) e 5 ( oito cursos ).
Um curso ficou Sem Conceito ( SC ).
Os Centros Universitários participaram com 38 cursos na região Sudeste, dos quais 15 obtiveram o conceito modal 3.
O conceito 1 foi atribuído a um curso, o conceito 2, a 12 cursos, e o conceito 4, a cinco cursos.
Nenhum curso recebeu o conceito 5.
Cinco cursos ficaram Sem Conceito ( SC ) nessa combinação de organização acadêmica e região.
Já as Faculdades foram representadas por 41 cursos na região Sudeste, dos quais três ficaram Sem Conceito ( SC ).
O conceito modal 3 foi atribuído a 22 cursos.
Aos demais cursos foram atribuídos os conceitos 1 ( três cursos ), 2 ( nove cursos ), 4 ( quatro cursos ).
Nenhum curso recebeu conceito 5.
Os CEFET/IF participaram com três cursos, todos receberam conceito 4.
Dos 60 cursos da região Sul, 39 eram mantidos por Universidades, para os quais o conceito modal foi 3, atribuído a 17 cursos.
O segundo conceito mais frequente foi 4, com 13 cursos.
O conceito 2 foi atribuído a quatro cursos, e o conceito 5 a três cursos.
Um curso recebeu conceito 1 e um curso ficou Sem Conceito ( SC ).
Dos 12 cursos participantes de Centros Universitários da região Sul, nove cursos foram avaliados com conceito 3 ( conceito modal ), dois cursos foram avaliados com conceito 2 e um curso avaliado com conceito 1.
Nenhum curso recebeu os demais conceitos ou ficou Sem Conceito.
Quanto aos nove cursos vinculados a Faculdades na região Sul, quatro receberam o conceito modal 3.
O conceito 2 foi atribuído a três cursos, e o conceito 4, a dois cursos.
Nenhum curso recebeu os conceitos 1 e 5 ou ficou Sem Conceito ( SC ).
Já os CEFET/IF não participaram na avaliação de cursos.
Na região Centro-Oeste, 16 dos 24 cursos oferecidos eram mantidos por Universidades.
Para este tipo de Organização Acadêmica, um curso ficou Sem Conceito ( SC ).
O conceito modal foi 4 ( cinco cursos ), seguido dos conceitos 3 ( quatro cursos ), 2 ( três cursos ), 5 ( dois cursos ) e 1 ( um curso ).
Os Centros Universitários da região Centro-Oeste contaram com apenas um curso, que foi avaliado com conceito 2.
As Faculdades participaram com 5 cursos e a todos foi atribuído o conceito 2.
Já os CEFET/IF participaram com dois cursos na região Centro-Oeste, aos quais foram atribuídos os conceitos 3 e 4.
No Gráfico 5.3, apresenta -se a distribuição do Conceito Enade dos cursos de Bacharelado, segundo a Organização Acadêmica da IES.
Nas Faculdades ( linha vermelha ) e nos CEFET/IF ( linha laranja ), todos os cursos têm conceito 4 ou abaixo disso, já que a poligonal cumulativa alcança 100 \% para o conceito 4.
Os cursos em CEFET/IF ( linha laranja ) apresentam uma poligonal mais à direita do que as demais, denotando uma melhor distribuição de Conceitos Enade.
No outro extremo, com distribuição muito semelhante entre si, encontram -se os cursos mantidos por Centros Universitários ( linha verde ) e Faculdades ( linha vermelha ), apresentando as respectivas poligonais mais à esquerda e, consequentemente, distribuição de Conceitos Enade com menores valores.
As linhas poligonais relativas às Faculdades ( em vermelho ) e aos CEFET/IF ( em laranja ) atingem 100 \% no conceito 4, já que nenhum curso dessas Organizações Acadêmicas recebeu conceito 5.

% matched lemmas: 
\end{textsample}
